% makes the page borders visible
% \usepackage{showframe}
% \usepackage{blindtext}


% \tolerance=1
\emergencystretch=\maxdimen
\hyphenpenalty=10000
\hbadness=10000



\usepackage{graphicx}
\usepackage{xcolor}
\usepackage{xstring,xifthen}
%%%%%%%%%%%%%%%%%%%%%%%%%%%%%%%%%%%%%%%%%%
%%%%% Template colors
%%%%%%%%%%%%%%%%%%%%%%%%%%%%%%%%%%%%%%%%%%
\definecolor{Prune}{RGB}{99,0,60} % l14-33 : couleurs de la charte graphique upsaclay
\definecolor{UPSgray1}{RGB}{49,62,72} % l14-33 : couleurs de la charte graphique upsaclay
\definecolor{UPSgray2}{RGB}{49,62,72} % l14-33 : couleurs de la charte graphique upsaclay
\definecolor{UPSgray3}{RGB}{49,62,72} % l14-33 : couleurs de la charte graphique upsaclay
\definecolor{UPSblue1}{RGB}{0,78,125} % l14-33 : couleurs de la charte graphique upsaclay
\definecolor{A1}{RGB}{99,0,60} % l14-33 : couleurs de la charte graphique upsaclay
\definecolor{B1}{RGB}{49,62,72}
\definecolor{C1}{RGB}{124,135,143}
\definecolor{D1}{RGB}{213,218,223}
\definecolor{A2}{RGB}{198,11,70}
\definecolor{B2}{RGB}{237,20,91}
\definecolor{C2}{RGB}{238,52,35}
\definecolor{D2}{RGB}{243,115,32}
\definecolor{A3}{RGB}{124,42,144}
\definecolor{B3}{RGB}{125,106,175}
\definecolor{C3}{RGB}{198,103,29}
\definecolor{D3}{RGB}{254,188,24}
\definecolor{A4}{RGB}{0,78,125}
\definecolor{B4}{RGB}{14,135,201}
\definecolor{C4}{RGB}{0,148,181}
\definecolor{D4}{RGB}{70,195,210}
\definecolor{A5}{RGB}{0,128,122}
\definecolor{B5}{RGB}{64,183,105}
\definecolor{C5}{RGB}{140,198,62}
\definecolor{D5}{RGB}{213,223,61}


\usepackage{mdframed}
\usepackage{titling}
\usepackage{emptypage}



\usepackage{tikz}
\usepackage{pgfplots}
\pgfplotsset{compat=1.16}
\usetikzlibrary{calc,angles,quotes}



\usepackage{chngcntr}
\counterwithout{equation}{chapter}
\counterwithout{figure}{chapter}
\counterwithout{table}{chapter}

\usepackage{perpage}
\MakePerPage{footnote}


\newlength{\oldparskip}
\newlength{\newparskip}
\setlength{\oldparskip}{\parskip}
\usepackage{parskip}
\setlength{\newparskip}{\parskip}

%%% Font
\usepackage{libertine}



\usepackage{scrextend} %% Forcer la 4ème  de couverture en page pair



\usepackage[absolute]{textpos}
\usepackage{colortbl}
\usepackage{array}





\usepackage[noindentafter,explicit]{titlesec}
\usepackage{titletoc}
\usepackage[chapter]{tocbibind}

\usepackage[twoside]{fancyhdr}
\pagestyle{fancy}



\usepackage{enumitem}
\usepackage{pifont}
\usepackage{fourier-orns}

%%% hyperref needs to be loaded at the perfectly right time for everything to work...
\usepackage[unicode=true,pdfusetitle,
 bookmarks=true,bookmarksnumbered=true,bookmarksopen=true,bookmarksopenlevel=5,
 breaklinks=false,pdfborder={0 0 0},pdfborderstyle={},backref=false,colorlinks=true]
 {hyperref}
 
\definecolor{titles}{HTML}{336084}
\hypersetup{linkcolor={titles},citecolor={A5},urlcolor={A3}}



\makeatletter
\def\@afterindentfalse{\let\if@afterindent\iffalse}
\@afterindentfalse





%%%%%%%%%%%%%%%%%%%%%%%%%%%%%%%%%%%%%%%%%%%%%%%%%%%%%%%%%


\def\isnum#1{%
  \if!\ifnum9<1#1!\fi%
     \emph{#1}%
  \else#1%
  \fi}
  
\newcommand*{\rom}[1]{\expandafter\@slowromancap\romannumeral #1@}



% \charcount 
% \spaces
% \wordchars
%%%%%%%%%%%%%%%%%%%%%%%%%%%%%%%%%%%%%%%%%%
%%%%% Numbering
%%%%%%%%%%%%%%%%%%%%%%%%%%%%%%%%%%%%%%%%%%
\renewcommand{\thechapter}{\Roman{chapter}}
\renewcommand{\thesection}{\thechapter-\arabic{section}}

%%%%%%%%%%%%%%%%%%%%%%%%%%%%%%%%%%%%%%%%%%
%%%%% Toc (titletoc package)
%%%%%%%%%%%%%%%%%%%%%%%%%%%%%%%%%%%%%%%%%%

\contentsmargin[1em]{2em}

\definecolor{tocvrules}{HTML}{8B8B8B}
\dottedcontents{}[]{}{}{}

\newcommand{\MainTocDottedRule}{%
{\titlerule*[8pt]{.}}%
}


%%%%%%%%%%%%%%%%%%%%%%%%%
%%% chapter toc content

\newcommand{\formatstarredchapter}[1]{{\normalsize #1}}
 
\newcommand{\formatcontentchapter}[1]{%
\parbox[b]{0pt}{%
\vspace{1pt}\makebox[0pt][r]{\huge \hyperpage{\thecontentspage}}%
}% end parbox with page
\hspace{6pt}%
{\color{tocvrules}\vrule width 1.5pt }% vrule between page and title
\hspace{6pt}%
\parbox[c]{0.7\textwidth}{%
%\vspace{2pt}%
\if!\ifnum9<1\thecontentslabel!\fi
\Large\textsc{\chaptertitlename} \thecontentslabel\vspace{1pt}\\%
\else%
\Large\textsc{\IfSubStr{ABCDEFGH}{\thecontentslabel}{Annexe}{\chaptertitlename}} \thecontentslabel\vspace{1pt}\\%
\fi%
\normalsize #1%
%\vspace{2pt}%
}% end parbox with title
\vspace{3pt}%
}% end \formatcontentchapter


\newcommand{\formatcontentchapterSimple}[1]{%
%\vspace{2pt}%
% \if!\ifnum9<1\thecontentslabel!\fi
% \Large\textsc{\chaptertitlename} \thecontentslabel\vspace{1pt}\\%
% \else%
\Large\textsc{\IfSubStr{ABCDEFGH}{\thecontentslabel}{Annexe}{\chaptertitlename}} \thecontentslabel\vspace{1pt}\\*%
% \fi%
\normalsize #1%
\vspace{3pt}%
}% end \formatcontentchapterSimple

\titlecontents{chapter}%
[0pt]% left
{\addvspace{1pc}%
% \IfStrEq{III}{\thecontentslabel}{\clearpage}{}%
}% above code
{\formatcontentchapterSimple}% for normal chapter
{\contentsmargin{0pt}%
\formatstarredchapter}% for starred chapter
{%
\StrLen{\thecontentslabel}[\mystringlen]%
    \ifthenelse{\mystringlen > 0}% test whether chapter is starred (by getting the length of the label)
    {%
        \hspace*{1ex}\MainTocDottedRule\contentspage[\textit{\thecontentspage}]%
    }%TRUE
    {%FALSE begin
        \hspace*{1ex}\MainTocDottedRule\contentspage[\textit{\thecontentspage}]%
    }%FALSE end
}% rule and page of all chapter
[%
% \needspace{1000pt}%
]%below spacing


%%%%%%%%%%%%%%%%%%%%%%%%%
%%% section toc content
\titlecontents{section}%
[3.8em] % ie, 1.5em (chapter) + 2.3em
{}%
{\contentslabel{2.3em}}%
{\hspace*{-2.3em}}%
{\hspace*{1ex}\MainTocDottedRule\contentspage[\textit{\thecontentspage}]}%
[\vspace{1pt}]%below spacing




%%%%%%%%%%%%%%%%%%%%%%%%%
%%% subsection toc content
\titlecontents{subsection}%
[7.1em] % ie, 1.5em (chapter) + 2.3em + 3.3em
{}%
{\contentslabel{3.3em}}%
{\hspace*{-2.3em}}%
{\hspace*{1ex}\MainTocDottedRule\contentspage[\textit{\thecontentspage}]}%
[\vspace{1pt}]%below spacing




%\renewcommand{\contentsname}{\protect\hypertarget{toctitle}{Table of contents}}


%%%%%%%%%%%%%%%%%%%%%%%%%%%%%%%%%%%%%%%%%%
%%%%% Partial Toc (titletoc package)
%%%%%%%%%%%%%%%%%%%%%%%%%%%%%%%%%%%%%%%%%%

\newcommand{\chapterPartialTocDottedRule}{%
{\titlerule*[8pt]{.}}%
}

\titlecontents{lsection}%
[3em] % ie, 1.5em (chapter) + 2.3em
{}%
{\contentslabel{2em}}%
{}%
{\hspace*{1ex}\chapterPartialTocDottedRule\contentspage[\textit{\thecontentspage}]}%
[\vspace{1pt}]%below spacing

% \makebox[2em][l]{\arabic{section}}
\titlecontents{lsubsection}%
[6em] % ie, 1.5em (chapter) + 2.3em + 3.3em
{}%
{\contentslabel{3em}}%
{}%
{\hspace*{1ex}\chapterPartialTocDottedRule\contentspage[\textit{\thecontentspage}]}%
[\vspace{1pt}]%below spacing


%%%%%%%%%%%%%%%%%%%%%%%%%
%%% partial toc

% \newcommand\PartialTocName{\contentsname}

\newcommand{\chapterPartialTocRule}{%
{\color{A4}\titlerule[2.5pt]}%
}

\def\chapterPartialTocDepth{1}

\newcommand{\chapterPartialToc}{%
{%
\setlength{\parskip}{\oldparskip}
% \vspace*{0.5pc}%
\vbox{\centering\Large \textit{Som\hspace{-0.7pt}maire du chapitre}}%
\vspace*{3pt}%
\chapterPartialTocRule
\vspace*{3pt}%
\printcontents{l}{1}[\chapterPartialTocDepth]{}}%
\vspace*{3pt}%
\chapterPartialTocRule
\vspace*{1pc}%
\setlength{\parskip}{\newparskip}
}



%%%%%%%%%%%%%%%%%%%%%%%%%%%%%%%%%%%%%%%%%%
%%%%% Titles customisation (titlesec package)
%%%%%%%%%%%%%%%%%%%%%%%%%%%%%%%%%%%%%%%%%%




%%%%%%%%%%%%%%%%%%%%%%%%%
%%% chapter format


\def\thickhrulefill{\leavevmode \leaders \hrule height 10pt depth -4pt \hfill \kern \z@}

% \newcommand{\chapterbreak}{\cleardoublepage}
\newcommand{\chapterbreak}{}

% \titleformat{\chapter}%
% [display]% shape
% {\huge\thispagestyle{empty}}% format
% {
% \filright\textsc{\chaptertitlename\ \thechapter}
% }% label
% {10pt}% sep between label and title
% {\filright\Huge\upshape\bfseries{#1}}%
% [\startcontents]% after code

\titleformat{\chapter}%
[display]% shape
{\huge\thispagestyle{empty}}% format
{% label
\centering \reset@font
\thickhrulefill\quad\scshape\chaptertitlename{} \thechapter\quad \thickhrulefill
}% label
{10pt}% sep between label and title
{\centering \vspace{10pt}%
% \hrule\vspace{2pt}%
% \Huge\upshape\bfseries{#1}%
\hrule
\vskip 2pt%
\Huge \bfseries #1
}%
[% after code
\vskip 5pt%
\hrule%
\startcontents]

\titlespacing*{\chapter}{0pt}{0pt}{30pt}


%\titleformat{\chapter}[hang]{\bfseries\Large\color{Prune}}{\thechapter\ -}{.1ex}
%{\vspace{0.1ex}
%}
%[\vspace{1ex}]
%\titlespacing{\chapter}{0pc}{0ex}{0.5pc}


\newcommand{\marginsecnumber}[1]{%
\makebox[0pt][r]{%
{#1}\hspace{8pt}%
}%
}
% \makebox[3em][l]{\thesubsection}
% \makebox[2em][r]{}
%%%%%%%%%%%%%%%%%%%%%%%%%
%%% section format
\titleformat{\section}
[hang]
{\bfseries\LARGE}
{\marginsecnumber{\arabic{section}}}
{0em}
{#1}
[]

\titlespacing*{\section}{0pt}{*3}{*2}
% \newcommand{\sectionbreak}{\clearpage}
%%%%%%%%%%%%%%%%%%%%%%%%%
%%% subsection format
\titleformat{\subsection}[hang]%shape
{\bfseries\Large}%format
{\marginsecnumber{\arabic{section}.\arabic{subsection}}}%label
{0em}%sep
{#1}%before code
[]%after code

\titlespacing*{\subsection}{0pt}{*3}{*2}

%%%%%%%%%%%%%%%%%%%%%%%%%
%%% subsubsection format
\titleformat{\subsubsection}[hang]%shape
{\bfseries\large}%format
{\marginsecnumber{\arabic{section}.\arabic{subsection}.\arabic{subsubsection}}}%label
{0em}%sep
{#1}%before code
[]%after code

\titlespacing*{\subsubsection}{0pt}{*3}{*2}

%%%%%%%%%%%%%%%%%%%%%%%%%
%%% paragraph format
% \titleformat{\paragraph}[runin]
% {\normalfont\normalsize\bfseries}{\theparagraph}{1em}{}[.]
\titleformat{\paragraph}[hang]%
{\bfseries\normalsize}%format
{\marginsecnumber{\thesubsubsection}}%label
{0em}%sep
{#1}%before code
[]%after code

\titlespacing*{\paragraph}{0pt}{*3}{*2}

%%%%%%%%%%%%%%%%%%%%%%%%%
%%% subparagraph format
\titleformat{\subparagraph}[runin]
{\normalfont\normalsize\bfseries}{\thesubparagraph}{1em}{#1}[.]

\renewcommand{\dotfill}{%
  \leavevmode\cleaders\hbox to 0.25em{\hss .\hss }\hfill\kern0pt }



%%%%%%%%%%%%%%%%%%%%%%%%%%%%%%%%%%%%%%%%%%
%%%%% Headers customisation (fancyhdr package)
%%%%%%%%%%%%%%%%%%%%%%%%%%%%%%%%%%%%%%%%%%

\fancypagestyle{plain}{}
\renewcommand{\chaptermark}[1]{\markboth{\emph{CHAPITRE \thechapter{} - #1}}{}}
\renewcommand{\sectionmark}[1]{\markright{\emph{Section \arabic{section}. #1}}}
\fancyhead[ER]{\leftmark}
\fancyhead[EL]{\thepage}

\fancyhead[OR]{\thepage}
\fancyhead[OL]{\rightmark}

\fancyfoot{}

% \renewcommand{\headrulewidth}{0pt} 

%%%%%%%%%%%%%%%%%%%%%%%%%%%%%%%%%%%%%%%%%%
%%%%% Biblio (biblatex)
%%%%%%%%%%%%%%%%%%%%%%%%%%%%%%%%%%%%%%%%%%
\AfterPackage{biblatex}{%
\renewcommand*{\bibfont}{\normalfont\footnotesize}%
}




%%%%%%%%%%%%%%%%%%%%%%%%%%%%%%%%%%%%%%%%%%
%%%%% list custom  (enumitem package) (just me having some fun)
%%%%%%%%%%%%%%%%%%%%%%%%%%%%%%%%%%%%%%%%%%

\setlist[itemize,1]{label={\color{Prune}{\scriptsize\ding{110}}}}
\setlist[itemize,2]{label={\color{Prune}{\scriptsize\ding{108}}}}
\setlist[itemize,2]{label={\color{Prune}{\scriptsize\ding{228}}}}

\setlist[itemize,1]{label={\color{Prune}{\scriptsize\decothreeright}}}
\setlist[itemize,2]{label={\color{Prune}{\scriptsize\decofourright}}}
% \setlist[itemize,2]{label={\color{Prune}{\scriptsize\ding{228}}}}

\usetikzlibrary{shadings}
\usetikzlibrary{shapes.geometric}

\def\sphericalbullet{%
\begin{tikzpicture}[scale=0.7]
  \shade[ball color = Prune!60,   opacity=0.9] (0,0) circle (5pt);
\end{tikzpicture}
}

\def\conicalbullet{%
\begin{tikzpicture}[scale=0.06,rotate=-90,transform shape]
    \shade[ball color = Prune!60,   opacity=0.9] (2,0) -- (0,6) -- (-2,0) arc (180:360:2cm and 0.5cm);
\end{tikzpicture}%
}


% \setlist[itemize,1]{label={{\sphericalbullet}}}
% \setlist[itemize,2]{label={{\conicalbullet}}}


\newcommand{\drawpolygon}[1]{%
\begin{tikzpicture}[]
\node[draw,minimum size=10pt,regular polygon,regular polygon sides=#1,inner sep=0] (a) {};
\end{tikzpicture}
}
% \setlist[itemize,1]{label={\drawpolygon{3}}}
% \setlist[itemize,2]{label={\drawpolygon{4}}}
% \setlist[itemize,3]{label={\drawpolygon{5}}}

\makeatother